% ========================================================
% BÀI TẬP TOÁN 5 - KẾT NỐI TRI THỨC VỚI CUỘC SỐNG
% Bài 25: Hình tam giác. Diện tích hình tam giác
% Định dạng: BTTeX Part III - Trắc nghiệm điền số ngắn
% ========================================================

% Câu 1
\begin{ex}
  Một hình tam giác có độ dài đáy là $12\text{ cm}$ và chiều cao là $8\text{ cm}$. Diện tích hình tam giác đó là bao nhiêu xăng-ti-mét vuông?
  \shortans{48}
  \loigiai{
    $(12 \times 8) : 2 = 48\text{ cm}^2$.
  }
\end{ex}

% Câu 2
\begin{ex}
  Một tam giác có đáy $15\text{ m}$ và chiều cao $6\text{ m}$. Diện tích tam giác đó là bao nhiêu mét vuông?
  \shortans{45}
  \loigiai{
    $(15 \times 6) : 2 = 45\text{ m}^2$.
  }
\end{ex}

% Câu 3
\begin{ex}
  Một tam giác vuông có hai cạnh góc vuông lần lượt là $9\text{ cm}$ và $12\text{ cm}$. Diện tích tam giác vuông đó là bao nhiêu xăng-ti-mét vuông?
  \shortans{54}
  \loigiai{
    $(9 \times 12) : 2 = 54\text{ cm}^2$.
  }
\end{ex}

% Câu 4
\begin{ex}
  Một tam giác vuông có diện tích $60\text{ cm}^2$, một cạnh góc vuông dài $10\text{ cm}$. Cạnh góc vuông còn lại dài bao nhiêu xăng-ti-mét?
  \shortans{12}
  \loigiai{
    Cạnh còn lại là: $(60 \times 2) : 10 = 12\text{ cm}$.
  }
\end{ex}

% Câu 5
\begin{ex}
  Một tam giác có diện tích $75\text{ m}^2$, chiều cao là $10\text{ m}$. Độ dài cạnh đáy tương ứng là bao nhiêu mét?
  \shortans{15}
  \loigiai{
    Đáy là: $(75 \times 2) : 10 = 15\text{ m}$.
  }
\end{ex}

% Câu 6
\begin{ex}
  Một tam giác có đáy dài $20\text{ cm}$, diện tích bằng $120\text{ cm}^2$. Chiều cao tương ứng là bao nhiêu xăng-ti-mét?
  \shortans{12}
  \loigiai{
    Chiều cao là: $(120 \times 2) : 20 = 12\text{ cm}$.
  }
\end{ex}

% Câu 7
\begin{ex}
  Nếu độ dài đáy của một hình tam giác tăng lên gấp đôi và chiều cao giữ nguyên thì diện tích của nó tăng lên gấp mấy lần?
  \shortans{2}
  \loigiai{
    Diện tích tỉ lệ thuận với độ dài đáy nên tăng 2 lần.
  }
\end{ex}

% Câu 8
\begin{ex}
  Nếu cả độ dài đáy và chiều cao của hình tam giác đều tăng lên gấp đôi thì diện tích của nó tăng lên gấp mấy lần?
  \shortans{4}
  \loigiai{
    Diện tích tăng $2 \times 2 = 4$ lần.
  }
\end{ex}

% Câu 9
\begin{ex}
  Một thửa ruộng hình tam giác có đáy dài $40\text{ m}$, chiều cao bằng $\frac{3}{4}$ đáy. Diện tích thửa ruộng là bao nhiêu mét vuông?
  \shortans{600}
  \loigiai{
    Chiều cao là $40 \times \frac{3}{4} = 30\text{ m}$. Diện tích là $(40 \times 30) : 2 = 600\text{ m}^2$.
  }
\end{ex}

% Câu 10
\begin{ex}
  Một bồn hoa hình tam giác có đáy $8\text{ m}$, chiều cao $5\text{ m}$. Người ta trồng hoa với chi phí $50\,000$ đồng mỗi mét vuông. Tổng chi phí trồng hoa là bao nhiêu triệu đồng?
  \shortans{1}
  \loigiai{
    Diện tích là $(8 \times 5) : 2 = 20\text{ m}^2$. Chi phí là $20 \times 50\,000 = 1\,000\,000$ đồng = 1 triệu đồng.
  }
\end{ex}

% Câu 11
\begin{ex}
  Một tam giác có đáy $14\text{ dm}$ và chiều cao $50\text{ cm}$. Diện tích tam giác đó là bao nhiêu đề-xi-mét vuông?
  \shortans{35}
  \loigiai{
    $50\text{ cm} = 5\text{ dm}$. Diện tích là $(14 \times 5) : 2 = 35\text{ dm}^2$.
  }
\end{ex}

% Câu 12
\begin{ex}
  Một tam giác có đáy $2{,}5\text{ m}$ và chiều cao $4\text{ m}$. Diện tích tam giác đó là bao nhiêu mét vuông?
  \shortans{5}
  \loigiai{
    $(2{,}5 \times 4) : 2 = 10 : 2 = 5\text{ m}^2$.
  }
\end{ex}

% Câu 13
\begin{ex}
  Tam giác $ABC$ có đáy $BC = 18\text{ cm}$. Kéo dài đáy $BC$ thêm $4\text{ cm}$ thì diện tích tăng thêm $24\text{ cm}^2$. Chiều cao của tam giác $ABC$ là bao nhiêu xăng-ti-mét?
  \shortans{12}
  \loigiai{
    Chiều cao là: $(24 \times 2) : 4 = 12\text{ cm}$.
  }
\end{ex}

% Câu 14
\begin{ex}
  Ở câu trên, diện tích tam giác $ABC$ ban đầu là bao nhiêu xăng-ti-mét vuông?
  \shortans{108}
  \loigiai{
    Diện tích là $(18 \times 12) : 2 = 108\text{ cm}^2$.
  }
\end{ex}

% Câu 15
\begin{ex}
  Một tấm biển quảng cáo hình tam giác có đáy $3\text{ m}$, chiều cao $2\text{ m}$. Sơn cả hai mặt tấm biển hết bao nhiêu mét vuông?
  \shortans{6}
  \loigiai{
    Diện tích một mặt là $(3 \times 2) : 2 = 3\text{ m}^2$. Cả hai mặt là $3 \times 2 = 6\text{ m}^2$.
  }
\end{ex}

% Câu 16
\begin{ex}
  Tam giác có tổng độ dài đáy và chiều cao là $28\text{ cm}$, đáy hơn chiều cao $4\text{ cm}$. Diện tích tam giác đó là bao nhiêu xăng-ti-mét vuông?
  \shortans{96}
  \loigiai{
    Đáy là $(28 + 4) : 2 = 16\text{ cm}$. Chiều cao là $12\text{ cm}$. Diện tích là $(16 \times 12) : 2 = 96\text{ cm}^2$.
  }
\end{ex}

% Câu 17
\begin{ex}
  Tam giác vuông có hai cạnh góc vuông hơn kém nhau $2\text{ cm}$ và tổng của chúng bằng $14\text{ cm}$. Diện tích tam giác vuông đó là bao nhiêu xăng-ti-mét vuông?
  \shortans{24}
  \loigiai{
    Hai cạnh là 8 cm và 6 cm. Diện tích là $(8 \times 6) : 2 = 24\text{ cm}^2$.
  }
\end{ex}

% Câu 18
\begin{ex}
  Một tam giác đều có chu vi $36\text{ cm}$. Độ dài mỗi cạnh là bao nhiêu xăng-ti-mét?
  \shortans{12}
  \loigiai{
    Cạnh tam giác đều là $36 : 3 = 12\text{ cm}$.
  }
\end{ex}

% Câu 19
\begin{ex}
  Một hình chữ nhật có diện tích $80\text{ cm}^2$. Nối một đường chéo chia hình chữ nhật thành hai hình tam giác bằng nhau. Diện tích mỗi tam giác là bao nhiêu xăng-ti-mét vuông?
  \shortans{40}
  \loigiai{
    $80 : 2 = 40\text{ cm}^2$.
  }
\end{ex}

% Câu 20
\begin{ex}
  Một lá cờ hình tam giác vuông có hai cạnh góc vuông là $20\text{ cm}$ và $30\text{ cm}$. Cần bao nhiêu xăng-ti-mét vuông vải để may $10$ lá cờ như thế?
  \shortans{3000}
  \loigiai{
    Diện tích 1 lá cờ: $(20 \times 30) : 2 = 300\text{ cm}^2$. 10 lá cờ cần $300 \times 10 = 3\,000\text{ cm}^2$.
  }
\end{ex}

% Câu 21
\begin{ex}
  Cho tam giác $ABC$ có diện tích $60\text{ cm}^2$. Điểm $M$ là trung điểm của cạnh $BC$. Diện tích tam giác $ABM$ là bao nhiêu xăng-ti-mét vuông?
  \shortans{30}
  \loigiai{
    Tam giác $ABM$ có chung chiều cao hạ từ $A$ và đáy bằng một nửa đáy $BC$, nên diện tích bằng $60 : 2 = 30\text{ cm}^2$.
  }
\end{ex}

% Câu 22
\begin{ex}
  Tam giác $ABC$ có diện tích $90\text{ cm}^2$. Điểm $N$ trên cạnh $AC$ sao cho $AN = \frac{1}{3} AC$. Diện tích tam giác $ABN$ là bao nhiêu xăng-ti-mét vuông?
  \shortans{30}
  \loigiai{
    Diện tích tam giác $ABN = \frac{1}{3}$ diện tích $ABC$, bằng $90 \times \frac{1}{3} = 30\text{ cm}^2$.
  }
\end{ex}

% Câu 23
\begin{ex}
  Một hình tam giác có đáy $16\text{ cm}$. Nếu giảm đáy đi $4\text{ cm}$ thì diện tích giảm đi $20\text{ cm}^2$. Chiều cao tương ứng của tam giác là bao nhiêu xăng-ti-mét?
  \shortans{10}
  \loigiai{
    Chiều cao là: $(20 \times 2) : 4 = 10\text{ cm}$.
  }
\end{ex}

% Câu 24
\begin{ex}
  Ở bài toán trên, diện tích tam giác lúc đầu là bao nhiêu xăng-ti-mét vuông?
  \shortans{80}
  \loigiai{
    Diện tích ban đầu là: $(16 \times 10) : 2 = 80\text{ cm}^2$.
  }
\end{ex}

% Câu 25
\begin{ex}
  Một tam giác vuông có diện tích $36\text{ cm}^2$ và hai cạnh góc vuông bằng nhau (tam giác vuông cân). Bình phương độ dài cạnh góc vuông là bao nhiêu?
  \shortans{72}
  \loigiai{
    $(a \times a) : 2 = 36 \implies a^2 = 72$.
  }
\end{ex}
